
\section*{Abstract}

Early fundus screening is an effective approach for ocular disease diagnosis. In this setting, traditional machine learning (ML), vision-language model (VLM), and convolutional neural network (CNN) approaches each offer distinct advantages. Considering feature modeling capacity, data requirements, and deployment cost, we adopt CNN as the main modeling backbone. However, real-world clinical applications still face several challenges, including false positives and false negatives caused by the co-occurrence of multiple diseases, degraded cross-institutional generalization performance, and privacy leakage risk. Motivated by these issues, we propose ARC, which captures disease co-occurrence dependencies through Transformer mechanisms to reduce false positives and false negatives. We further design DFE, which learns domain-invariant features from a single source domain when multi-domain collaborative training is restricted. With CAM-guided training, DFE improves generalization to unseen domains, while CAM visualizes model-attended regions and enhances interpretability. In addition, AdaGDP is integrated into model training to preserve privacy through dynamic gradient noise injection while maintaining model performance under moderate privacy preservation. Experiments confirm the effectiveness of each module through comprehensive ablation studies and demonstrate performance improvements on three datasets under both in-domain and cross-domain generalization settings, with controllable performance degradation under privacy preservation. Overall, we establish a privacy-preserving interpretable single domain generalization framework for multi-label ocular disease recognition, improve cross-domain generalization performance for multi-disease recognition, and provide insights into ocular disease diagnosis under privacy-preserving scenarios.

\noindent\textbf{Keywords:}\textbf{ Domain Generalization; Multi-Label Ocular Disease Recognition; Data Privacy; Interpretability}
